\chapter{Data analysis techniques}
\label{ch:methods}

As a result of the quantum-mechanical nature of the collisions and the interaction of their products with the detectors, the observations resulting from a particular interaction are fundamentally probabilistic. 
%
Therefore, the approach to data analysis and the conclusions drawn from the data must be framed in statistical terms, which include not only low-level tasks, such as the reconstruction and identification of particles and the measurement of their energies and momenta, but also high-level tasks, such as searches for new particles and measurements.
%
The high dimensionality and large volume of LHC data pose a problem because the statistical model over the high-dimensional space of their experimental data is not known explicitly in terms of an equation that can be evaluated. 
%
Instead, one typically has access to large samples of simulated data generated by stochastic simulation programs that model the physics of particle interactions on various scales. 
%
If the data were fairly low dimensional (d< 5), the problem of estimating the unknown statistical model from the simulated samples would not be difficult.
%
Histograms or kernel-based density estimates provide reasonable estimates in low-dimensional spaces. 
%
These strategies, however, suffer from the curse of dimensionality. 
%
In a single dimension, \(N\) samples may be required to describe the source probability density function. 
%
In \(d\) dimensions, the number of samples required increases to the power of the data’s dimensionality: \(\mathcal{O}(N^d)\). 
%
The consequence is that any dimensionality greater than five or ten requires impractical or impossible computational resources, regardless of the speed of the sample generator.

%In parallel to the relatively recent rise of ML techniques in industrial applications, scientists have increasingly become interested in the potential of ML for fundamental research, and HEP is no exception.
%
%Traditionally, HEP physicists have approached this problem by reducing the dimensionality of the data through a series of steps that operate both on individual collision events and on collections of events. For an individual event, reconstruction algorithms process the raw sensor data into low-level objects such as calorimeter clusters and tracks. 
%%
%From these low-level components, the algorithms attempt to estimate the energy, momentum, and identity of individual particles.
%% 
%Next, from these reconstructed objects, event-level summaries are constructed. 
%%
%Event selection algorithms then select subsets of the collision data for further analysis on the basis of the information associated with individual events. 
%%
%Traditionally, the reconstruction and event selection operations are based on specific, engineered features in the data. 
\section{Machine Learning Basics}
Fortunately, many of the tasks encountered in HEP can be naturally reformulated as machine learning problems. 
%
Typically, the problems are formulated in terms of a search for some function \(f:X \to Y\), from the space of the observed data \(X\) to a low-dimensional space of a desired target label \(Y\), which optimizes some metric of our choosing. 
%
This metric is called a loss function and is written as \(L[ \mathbf{y}, f (\mathbf{x})]\).
%

Ideally, a learning algorithm would find the function that optimizes \(L\) over all possible values
of \((\mathbf{x}, \mathbf{y})\), but this is intractable owing to the curse of dimensionality and an infinite number of functions to choose from. 
%
Instead, in supervised learning, one has labeled training data \(\{\mathbf{x}_i , \mathbf{y}_i\}^N_{i=1}\) sampled from \(p(\mathbf{x} , \mathbf{y})\). 
%
Furthermore, the function space is restricted to a model—a highly flexible family of functions \(f_\phi(\mathbf{x})\) parameterized by \(\phi\). 
%
In this case, the algorithms minimize the loss function directly with respect to the model parameters \(\phi\). 
%
Neural networks, support vector machines, and decision trees are examples of types of models commonly used in machine learning. 
%
These models often have a large number of parameters, and in the case of neural networks, finding the optimal \(f_\phi\) can be a difficult problem
%

An essential goal in machine learning is generalization—the ability of the model to perform well on data that were not used in training. 
%
Failure in this task is called overtraining or overfitting. 
%
A vast array of techniques exist to avoid overtraining, all of which can be considered forms of regularization. 
%
Regularization techniques such as dropout were key to advances in image recognition with deep learning

\section{Jet Classification}


\section{Unfolding}

\section{Uncertainties in ML predictions}




